\documentclass[11pt]{article}
\usepackage[a4paper,margin=1in]{geometry}
\usepackage{amsmath,amssymb}
\usepackage{microtype}
\usepackage[hidelinks]{hyperref}

\title{\textbf{The missing \(\lambda>2\) injectivity range is answered}\\
\large Literature identification for arXiv:2201.07837}
\author{Automated mathematics run \texttt{fa\_banach\_001}\\
Lane 14, model GPT5.6}
\date{9 August 2026}

\begin{document}
\maketitle

\noindent\textbf{Status: literature\_already\_answered.}

\section*{Source question}

Kania and Lewicki's companion paper
\emph{A forgotten theorem of Pe{\l}czy\'nski:
\((\lambda+)\)-injective spaces need not be \(\lambda\)-injective - the case
\(\lambda\in(1,2]\)} (arXiv:2201.07837) proves the result for
\(1<\lambda\le2\). In the Introduction, immediately after Theorem A, it
states:
\begin{quote}
\emph{Furthermore, the complementary case of
\(\lambda\in(2,\infty)\) remains a mystery to us.}
\end{quote}
Thus the exact missing assertion is the existence, for every \(\lambda>2\),
of a \((\lambda^+)\)-injective Banach space which is not
\(\lambda\)-injective.

\section*{Answering theorem}

The same authors' later paper \emph{\((\lambda^+)\)-injective Banach spaces}
(arXiv:2603.09710, 2026) explicitly says in its abstract and Introduction
that it resolves the complementary range left open in the companion paper.
Its Theorem A states:
\begin{quote}
\emph{For every \(\lambda>1\) there exists a Banach space that is
\((\lambda^+)\)-injective but not \(\lambda\)-injective.}
\end{quote}
This contains the entire missing interval \((2,\infty)\), so it is an exact
affirmative answer.

\section*{Identification}

The answering paper introduces the zero-sum subspace
\[
 \Sigma_N(Y)=\left\{(y_1,\ldots,y_N)\in Y_\infty^N:
 \sum_{j=1}^N y_j=0\right\}.
\]
Its multiplication lemma shows that the relative projection constant is
multiplied by
\[
 \mu_N=2-\frac{2}{N},
\]
and that failure of attainment is preserved. Iterating this construction
reduces any prescribed \(\lambda>2\) to a base constant in \((1,2]\)
provided by arXiv:2201.07837.

\section*{Scope and provenance}

The supporting authors explicitly know they are completing their earlier
open range, so the appropriate status is
\texttt{literature\_already\_answered}, not a new proof produced by this
run. Together the two papers establish Pe{\l}czy\'nski's claimed theorem for
all \(\lambda>1\). The complex endpoint \(\lambda=1\) remains outside this
answer.

\begin{thebibliography}{2}
\bibitem{KL2023}
T. Kania and G. Lewicki,
\emph{A forgotten theorem of Pe{\l}czy\'nski:
\((\lambda+)\)-injective spaces need not be \(\lambda\)-injective - the case
\(\lambda\in(1,2]\)}, Studia Math. (2023), arXiv:2201.07837.

\bibitem{KL2026}
T. Kania and G. Lewicki,
\emph{\((\lambda^+)\)-injective Banach spaces},
arXiv:2603.09710 (2026).
\end{thebibliography}

\end{document}

